%!TEX root = TQFTnotes.tex
%%%%%%%%%%%%%%%%%%%%%%%%%%%%%%
\chapter{Dijkgraaf--Witten theory}\label{c:DW}
In this chapter we discuss the fundamental example of a TQFT:
the gauge theory with finite gauge group,  also known as the Dijkgraaf--Witten
model. This theory was introduced by Dijkgraaf in his thesis
\cite[Section 8.5]{dijkgraaf-thesis} (see also \cite[Section 6]{DW}). In
these works, the focus was mostly on the theory in dimension 3, where it can be
considered as a baby version of Chern--Simons theory, but it works the same way
in any dimension. We will present the Dijkgraaf--Witten model following
\cite{freed-quinn}, \cite{teleman-five-lectures}.

Throughout this chapter, we fix a finite group $G$; for simplicity, we will
use the ground field $\kk=\C$.

Informally, for a closed $n$-manifold $M$, $Z_G(M)$ should count $G$-covers
(or $G$-bundles) on $M$:
\begin{equation}\label{e:dw-informal}
  Z_G(M)=\text{``number of $G$-bundles $P\to M$''}.
\end{equation}
(The reason why this is in quotes will be discussed shortly.) For people
not familiar with the notion of $G$-bundles, we give a review of the main
facts.

%%%%%%%%%%%%%%%%%%%%%%%%%%%%%%%%%%%%%%
\section{Principal $G$-bundles}\label{s:principal-bundles}

\begin{definition}\index{G-bundle@$G$-bundle}
  A $G$-bundle over a manifold $M$ is a finite covering $\pi\colon P\to M$
  together with a right action of $G$ on $P$ which preserves $\pi$ and
  such that the action of $G$ on each fiber $P_x=\pi^{-1}(x)$ is free and
  transitive.
\end{definition}
The reason for choosing a right action instead of a left one is purely
technical; one could also use left action and get an equivalent definition.

We define an isomorphism of $G$-bundles (for the same $G$ and $M$) in the
obvious way. We denote by $\Bun_G(M)$ the category of $G$-bundles on $M$, with
morphisms being isomorphisms of $G$-bundles. This category is a
\emph{groupoid}\index{groupoid}, i.e., a category in which all morphisms are
invertible. In particular, for a $G$-bundle $P$ we denote by $\Aut(P)$ the group
of automorphisms of $P$.

Since $G$ is discrete, any $G$-bundle comes with a flat connection: given a
path $\ga$ in $M$ connecting points $x_0, x_1\in M$, and a choice of lift
$\tilde x_0\in P_{x_0}$, there is a unique way to lift $\ga$ to a path in $P$.
This gives a holonomy map $\rho_\ga \colon P_{x_0} \to P_{x_1}$ which commutes
with the right action of $G$. In particular, this gives us the monodromy map
\begin{equation}\label{e:monodromy}
  \rho \colon \pi_1(M,x)\to \Aut(P_x)
\end{equation}
where $\Aut(P_x)$ is the group of all bijections $P_x\to P_x$ which commute with
the right action of $G$.

If we choose a point $\tilde x\in P_x$, we can identify $P_x$ with $G$;
this gives an identification $\Aut(P_x)\isoto G$ (since any bijection $G\to G$
which commutes with the right action of $G$ is given by left multiplication
by some element of $G$). Thus, a pair $(P, \tilde x)$ gives rise to a
monodromy map $\rho \colon \pi_1(M,x)\to G$. It is easy to see that
conversely, for connected $M$, any group homomorphism $\pi_1(M,x)\to G$
defines a $G$-bundle $P$ together with a choice of a base point $\tilde x
\in P_x$. In other words, we have the following lemma.

\begin{lemma}\label{l:monodromy1}
  Let $M$ be a connected manifold; fix a choice of base point $x\in M$.
  Then the set of isomorphism classes of pairs $(P, \tilde x)$,
  $\tilde x\in P_x$, is in bijection with
  $$
  \Hom(\pi_1(M,x),G).
  $$
  Moreover, such pairs have no non-trivial automorphisms: if $(P, \tilde x)$
  and $(P', \tilde x')$ are two such pairs, then the isomorphism between them,
  if it exists, is unique.
\end{lemma}

We will refer to a pair $(P, \tilde x)$ as in the lemma as a
\emph{based $G$-bundle}\index{G-bundle@$G$-bundle!based}. One can restate
this lemma as follows: the groupoid $\tBun_G$
of based $G$-bundles is equivalent as a groupoid to the set
$\Hom(\pi_1(M,x), G)$ (considered as a trivial groupoid: the only morphisms
are identity morphisms).

If we change the choice of $\tilde x$, replacing it by $\tilde x g$ (for
the same $P$), then it is easy to see that the monodromy map
\eqref{e:monodromy} is conjugated by $g^{-1}$. Thus, we get the following
result.

\begin{lemma}\label{l:monodromy2}
  Let $M$ be connected. Then the set of isomorphism classes of objects
  in $\Bun_G(M)$ is in bijection with
  $$
    \Hom(\pi_1(M,x),G)/G
  $$
  where $G$ acts by conjugation. Moreover, for $\rho \in \Hom(\pi_1(M,x),G)$,
  the automorphism group of the corresponding bundle $P$ is
  \begin{equation}\label{e:aut-bundle-monodromy}
    \Aut(P) = \{g\in G\st g\rho g^{-1} = \rho\}.
  \end{equation}
\end{lemma}
One can restate it as an equivalence of groupoids.
\begin{definition}
  Let $X$ be a set with an action of a group $G$. We define the
  groupoid $X\q G$ as the category whose objects are points $x\in X$ and
  $\Hom(x_0, x_1)=\{g\in G\st g(x_0)=x_1\}$.
\end{definition}
Then \leref{l:monodromy2} says that for a connected $M$, the groupoid
$\Bun_G(M)$ is equivalent to $\Hom(\pi_1(M,x),G)\q G$.


\begin{example}\label{x:bundles_s1}
  Let $M=S^1$. Then $\tBun_G(S^1)\simeq G$, and $\Bun_G(S^1)\simeq G\q G$ (in
  the second identity, $G$ acts on $G$ by conjugation). In particular, there is
  a bijection between the set of isomorphism classes of $\Bun_G(S^1)$ and the
  set of conjugacy classes in $G$; if we denote by $P_g$ the bundle
  corresponding to $g\in G$, then $\Aut(P_g)$ is the centralizer of $g$:
  $\Aut(P_g)\simeq Z_g=\{h\in G\st hgh^{-1}=g\}$.
\end{example}

%%%%%%%%%%%%%%%%%%%%%%%%%%%%%%%%%%%%%%
\section{Definition of DW theory}\label{s:DW-definition}
In this section, we define the Dijkgraaf--Witten (DW) TQFT; this definition
works in any dimension $n$.

Naively, we want to define, for a closed $(n-1)$-manifold $N$, the vector space
$$
  Z_G(N)=\text{vector space generated by isomorphism classes of
    $G$-bundles on $N$}
$$
and for a cobordism $N_0\cob{M}N_1$, we define
$$
  Z_G(M)([P_0])=\sum_{P} [P|_{N_1}],\qquad P_0\in \Bun_G(N_0)
$$
where the sum is over all isomorphism classes of bundles $P$ on $M$ such that
$P|_{N_0}=P_0$. The gluing axiom follows immediately.

This, however, needs further work. First, the requirement $P|_{N_0}=P_0$
is meaningless: two objects in a category (in this case, two bundles) are never
equal; instead, we should talk about them being isomorphic. But even if $M$ is
closed, so we do not have this problem, just counting isomorphism classes gives
the wrong answer, as we completely ignore the fact that some of these bundles
have non-trivial automorphisms.

To begin with, let us explain how one counts $G$-bundles on a connected
manifold $M$. In this case, we can choose a base point $x\in M$; then for any
bundle $P$ on $M$, we have $|G|$ possible choices of lifts $\tilde x\in P_x$.
Thus, we make the following definition
\begin{equation}\label{e:count-bung}
  |\Bun_G(M)|=\frac{1}{|G|} |\tBun_G(M)|
\end{equation}
where we denoted by $|\tBun_G(M)|$ the number of isomorphism classes of
$(P,\tilde x)\in \tBun_G(M)$. Note that there are no non-trivial
automorphisms in $\tBun_G$: this groupoid is equivalent to a set.

\begin{lemma}\label{l:count_aut}
  Let $M$ be connected, and let $|\Bun_G(M)|$ be defined by
  \eqref{e:count-bung}. Then
  $$
    |\Bun_G(M)|=\sum_{[P]} \frac{1}{|\Aut(P)|}
  $$
  where the sum is taken over isomorphism classes in $\Bun_G(M)$.
\end{lemma}
In other words, in this count every isomorphism class is taken with
the weight equal to $\frac{1}{|\Aut(P)|}$.

This lemma immediately follows from the following elementary observation,
whose proof we leave to the reader: if $X$ is a set with an action of a
finite group $G$, then
$$
  \frac{|X|}{|G|}=\sum_{[x]}\frac{1}{|\Stab(x)|}
$$
where $\Stab(x)=\{g\in G\st g(x)=x\}=\Aut_{X\q G}(x)$.
\begin{example}\label{x:bun-circle}
  Let $M=S^1$. Then $\tBun_G(M)\simeq G$, so $|\Bun_G(M)|=\frac{1}{|G|}
  |\tBun_G(M)|=1$.
  For an element $g\in G$, the automorphism group of the corresponding bundle
  is the centralizer of $g$: $\Aut(P_g)=Z_g$; thus,
  $$
  \frac{1}{|\Aut(P_g)|}= \frac{|C_g|}{|G|}
  $$
  where $C_g\subset G$ is the conjugacy class containing $g$.
\end{example}

More generally, we see that for a connected manifold $M$, we have
\begin{equation}\label{e:dw-closed}
  Z_G(M)= |\Bun_G(M)|= \frac{|\Hom(\pi_1(M), G)|}{|G|}.
\end{equation}


\leref{l:count_aut} suggests the following definition, which now
can be used for any groupoid (in particular, for $\Bun_G(M)$ for general $M$).
\begin{definition}\label{d:count}\index{groupoid!cardinality}
  Let $\calC$ be a groupoid such that the set of isomorphism classes is finite,
  and for any $x\in \calC$, the group
  $\Aut(x)$ is finite. Then we define
  $$
  |\calC|=\sum_{[x]} \frac{1}{|\Aut(x)|}
  $$
  where the sum is taken over representatives of isomorphism classes in $\calC$.
\end{definition}
\begin{example}\label{x:action-groupoid}
  Let $X$ be a finite set with an action of a finite group $G$.
  Then $|X\q G|=|X|/|G|$. Note that we are not assuming that the
  action of $G$ is free.
\end{example}
The above formulas for counting elements in ``sets with automorphisms'' (aka
groupoids) have a long history; in particular, they have been extensively used
in physics, when discussing Feynman diagrams. For a nice discussion of this see
\cite{baez-dolan-finite-sets}. In particular, we can count how many finite sets
there are --- try doing this yourself!

See also the blog post \cite{baez-counting} by John Baez.


Similar ideas can also be used, e.g., to define the Euler characteristic of
an orbifold (see \cite[Section 2.4]{caramello}). Note: there are several
different definitions of the orbifold Euler characteristic.

Now that we have discussed the proper way of counting, we can return to the
definition of DW theory.

\begin{theorem}\label{t:DW}
  Let $G$ be a finite group. Then the formulas below define an
  $n$-dimensional TQFT, called the Dijkgraaf--Witten TQFT.

  For a closed $(n-1)$-manifold $N$, we define
  $$
    Z_G(N)=\text{vector space generated by isomorphism classes of
                  $G$-bundles on $N$}
  $$

  For a cobordism $N_0\cob{M}N_1$, we define the linear map
  $Z_G(M)\colon Z_G(N_0)\to Z_G(N_1)$ by
  $$
    Z_G(M)([P_0])=\sum_{(P, \ph)} \frac{1}{|\Aut(P,\ph)|} [P|_{N_1}]
  $$
  where the sum is over all isomorphism classes of pairs $(P, \ph)$, where
  $P$ is a $G$-bundle on $M$, $\ph$ is an isomorphism of
  bundles $P_0\to P|_{N_0}$, and $\Aut(P, \ph)$ is the group of
  automorphisms of $P$ which commute with $\ph$.
\end{theorem}

Note that if $M$ is connected and $N_0$ is non-empty, there are
no non-trivial automorphisms of $(P,\ph)$.

Equivalently, one can define %FIXME: verify!!
$$
  Z_G(M)([P_0]) = \sum_{[P_1]} \Bigl(
                                    \sum_{[P]} \frac{1}{|\Aut_0(P)|}
                                \Bigr) [P_1]
$$
where the outer sum is over all isomorphism classes of bundles on $N_1$, the
inner sum is over all isomorphism classes of bundles on $M$ such that
$P|_{N_0}\simeq P_0$, $P|_{N_1}\simeq P_1$, and $\Aut_0(P)$ denotes the group of
automorphisms of $P$ which are the identity on $N_0$. Some authors instead use
automorphisms that are identity on $N_1$; these two definitions are equivalent
after conjugation by the operator
$Z_G(N)\to Z_G(N)\colon [P]\mapsto \frac{[P]}{|\Aut(P)|}$.

\begin{example}\label{x:dw-invariants}\par\noindent
  \begin{enumerate}
    \item Let $M=S^n$ (we assume $n\ge 2$). Then any $G$-bundle on $M$ is
      trivial, and $\Aut(P)=G$, so $Z_G(S^n)=\frac{1}{|G|}$.
    \item Let $M=N\times I$ be the cylinder over $N$.
      Then for a given bundle $P_0$ on $N$,
      there is a unique bundle $P$ on $M$ extending it (more precisely, the pair
      $(P, \ph)$ is unique up to unique isomorphism); thus, $Z_G(M)=\id$ as
      expected.
    \item Let $M$ be the cylinder over $N$ considered as cobordism $N\sqcup
      \ov{N}\to \varnothing$. For a bundle $P$ on $N$,
      denote by $\bar P$ the same bundle considered as a bundle
      on $\ov{N}$. Then
      $$
      Z_G(P_1, \ov{P_2})=
      \begin{cases}
        |\Aut(P_1)| &\quad P_1\simeq P_2\\
        0 & \quad\text{otherwise}
      \end{cases}
      $$
      Thus, the isomorphism classes of bundles on $N$ form an orthogonal but not
      orthonormal basis with respect to the pairing given by this cobordism.
  \end{enumerate}
\end{example}

Let us now consider the special case $n=2$. In this case, every element $g\in G$
gives a bundle $P_g$ on $S^1$ with monodromy $g$, and $P_g\simeq P_{g'}$ iff $g,
g'$ are conjugate in $G$ (we will denote it $g\sim g'$); in particular, $P_1$ is
the trivial bundle. We will denote the corresponding element of $A=Z_G(S^1)$ by
$[P_g]$. By general theory (\thref{t:2d-classification2}), $A$ must have the
structure of a commutative Frobenius algebra. The following lemma gives an
explicit description of this algebra structure.

\begin{lemma}\label{l:dw-frobenius-structure}
  The unit, the counit, and the multiplication in $A$ are given by
  \begin{equation}\label{e:dw-frob-structure}
    \begin{aligned}
      1 &= \frac{1}{|G|} [P_1]\\
      \eps([P_g]) &= \de_{g,1}\\
      [P_{g_1}][P_{g_2}]&=\sum_{h\in G} [P_{g_1 hg_2h^{-1}}].
    \end{aligned}
  \end{equation}
\end{lemma}
The proof of this lemma is given by an explicit computation. For example, the
unit is given by $Z_G(\tinycap)$; since every bundle on a disk is trivial, and
the automorphism group of the trivial bundle is $G$, the definition of $Z_G(M)$
given in \thref{t:DW} becomes $\frac{1}{|G|} [P_1]$. The other cases are given
by similar explicit computations.

\begin{corollary}
  Let $\C[G]$ be the group algebra of $G$.
  Define the linear map $A\to \C[G]$ by
  $$
    [P_g]\mapsto \sum_{h \in G} hgh^{-1}.
  $$
  Then this map is an algebra isomorphism $A\isoto \C[G]^G=z(\C[G])$
  \textup{(}here, $z$ stands for the center of the algebra\textup{)}, which
  identifies $\eps$ with the map
  \begin{equation}\label{e:dw-eps}
    \eps(g) = \frac{1}{|G|}\de_{g, 1}= \frac{1}{|G|^2}\tr_{R}(g),
  \end{equation}
  where $R=\C[G]$ is the regular representation.
\end{corollary}

We can describe the algebra $A=\C[G]^G$ in the form of
\thref{t:frobenius_commutative}. Namely, let $\{V_i\}$ be the set of
representatives of isomorphism classes of irreducible representations of $G$;
recall that by Maschke's theorem, any complex finite-dimensional representation
of $G$ is isomorphic to a direct sum of copies of $V_i$. By a well-known result
of representation theory of finite groups, $\C[G]\simeq \bigoplus_i \End(V_i)$,
which implies
\begin{equation}\label{e:dw-projectors}
  A=\C[G]^G=\bigoplus_i \C e_i,
\end{equation}
where $e_i$ are the projectors onto irreducible representations:
$e_i|_{V_j}=\de_{ij}\id_{V_i}$; thus, $e_i e_j = \de_{ij}e_i$.

\begin{remark}\label{r:characters}
  It follows from orthogonality of characters that
  $$
    e_i =\frac{\dim V_i}{|G|} \sum_{g\in G}\chi_i(g^{-1})g,
  $$
  where $\chi_i(g)=\tr_{V_i}(g)$ is the character of $V_i$. Thus, characters of
  irreducible representations, considered as elements in $\C[G]^G$, are (up to
  scalars) exactly the orthogonal idempotents. In fact, this is how characters
  were first introduced in \cite{frobenius-characters}: in modern terminology,
  they were defined as elements of the center of the group algebra such that
  $\chi_i\chi_j = \de_{ij}c_i$. The statement that they are traces in
  irreducible representations (and indeed the notion of irreducible
  representation itself) appeared later.
\end{remark}

Since every $V_i$ appears in the regular representation with
multiplicity $d_i=\dim V_i$, \eqref{e:dw-eps} gives
$$
  \eps(e_i) = \la_i  = \frac{\dim^2 V_i}{|G|^2}.
$$
We can use this to compute $Z_G(\Si_g)$, where $\Si_g$ is a closed
surface of genus $g$. Indeed, by \leref{l:semisimple_theory}, we have
$$
  Z_G(\Si_g)= \sum_i \la_i^{1-g} = |G|^{2g-2} \sum_i (\dim V_i)^{2-2g}
$$
where the sum is over all isomorphism classes of irreducible representations.
On the other hand, by \eqref{e:dw-closed},
$$
  Z_G(\Si_g) =\frac{|\Hom(\pi_1(\Si_g), G)|}{|G|}.
$$
The structure of the fundamental group $\pi_1(\Si_g)$ is well-known. For $g=0$
it is trivial, and for $g\ge 1$ it is generated by elements $a_i, b_i$ with a
single relation $\prod_i [a_i,b_i]=1$, where $[a_i,b_i] =
a_ib_ia_i^{-1}b_i^{-1}$ is the group commutator. Thus, for $g\ge 1$,
we  we get the following identity
\begin{equation}\label{e:dw-genusg}
  |G|^{2g-1} \sum_i (\dim V_i)^{2-2g}
  =\bigl|
    \{a_1, \dots, a_g, b_1, \dots, b_g\in G\st \prod_i [a_i,b_i]=1\}
    \bigr|.
\end{equation}

In particular, for $g=1$ \eqref{e:dw-genusg} gives
$$
  |\{g,h\in G\st gh=hg\}|=|G|\cdot (\text{number of irreps of $G$})
$$
which can also be proved directly using representation theory methods without
any reference to TQFT. This formula was first obtained by Frobenius in 1896,
see \cite{frobenius-characters}.

%%%%%%%%%%%%%%%%%%%%%%%%
\section{Path integral interpretation of DW theory}\label{s:DW-path-integral}
One can rewrite our definition of DW theory so that it looks similar to the
original description of QFT in terms of path integrals. Namely, let $BG$
be the classifying space of the group $G$; this is a topological space which
is obtained by taking a contractible space $EG$ on which $G$ acts freely and
then defining $BG=EG/G$. (Explicit constructions of $EG$ and $BG$ are
well-known; note however that typically $BG$ is infinite-dimensional,
even for a finite group $G$).

We have the universal $G$-bundle $EG\to BG$; thus, any continuous map
$\ga\colon M\to BG$ gives rise to the pullback bundle
$P_\ga = \ga^*(EG)$ over $M$. It is well known that any $G$-bundle can be
obtained in this way; more precisely, we have a bijection
$$
\text{isomorphism classes in }\Bun_G(M)= [M, BG]
$$
where $[M,BG]$ is the set of homotopy classes of
maps $M\to BG$.

Thus, the definition of DW invariant of a closed $n$-manifold can
be rewritten as follows
$$
  Z_G(M)=\int_{[M, BG]} D\ga
$$
where the integral is over the set of homotopy classes of maps (which is a
finite set) and the measure $D\ga$ is defined by the condition that each
homotopy class $[\ga]$, considered as a point in the finite set $[M, BG]$, has
measure $\frac{1}{|\Aut(P_\ga)|}$.

This also suggests that we can define a ``twisted'' version of DW theory.
Assume from now on that $M$ is a closed oriented $n$-manifold, with
fundamental class $[M]\in H_n(M, \Z)$. Let
$\om\in H^n(BG, U(1))$ be a cohomology class; then $\ga$ allows us to
define the pullback class $\ga^*(\om)\in H^n(M, U(1))$, and pairing it with
$[M]$ gives a number. Since $\om$ is a cohomology class, the pullback
$\ga^*(\om)$ depends only on the homotopy class $[\ga]$, so the pairing
$\<\ga^* \om, [M]\>$ is a well-defined function on $[M, BG]$. Thus, we define
$$
  Z_G^\om (M)= \int_{[M, BG]}\<\ga^* \om, [M]\> D\ga.
$$
One should think of $\<\ga^* \om, [M]\>$ as an analogue of the term
$e^{iS(\phi)}$ in the path integral.

Note that the short exact sequence of abelian groups $0\to \Z\to \R \to \R/\Z\to
0$ gives rise to a long exact sequence of cohomology of $BG$; since it is well
known that for a finite group and $n\ge 1$, $H^n(BG, \R)=H^n(G,\R)=0$ (here
$H^n(G)$ is the group cohomology), we get an isomorphism
$$
  H^n(BG, U(1))=H^{n+1}(BG, \Z).
$$
Thus, the ``twist'' $\om$ can also be thought of as an element
in $H^{n+1}(BG,\Z)$.

\clearpage
